This section defines the case-study generator, constraints, cohorts, and the SEME endpoint families (Signed concordance, Effect-vector fidelity, Magnitude transfer, Encoding retention). Paired post-dose EEG is not part of the design. Configuration keys, scripts, and artifact paths are listed in Supplementary Material S1.

\subsection{Architecture and pharmacodynamic penalties}

The twin is a pharmacodynamically constrained CVAE (\texttt{constrained\_mlp}): MLP encoders, no graph EEG encoder, no cross-attention, constrained loss enabled. Input EEG is a packed 2185-D vector for 19 standard 10-20 channels (PSD \([0,380)\), band powers \([380,475)\), interleaved coherence/PLV connectivity \([475,2185)\)). \(L_{\mathrm{band}}\) uses delta through beta slices; \(L_{\mathrm{conn}}\) uses window \([475,1330)\), which mixes early-band coherence and PLV under the interleaved flatten (S1). Layer widths and parameter count are in Table~\ref{tab:architecture}. Fig.~\ref{fig:architecture} summarizes data flow.

\begin{figure*}[!t]
\centering
\includegraphics[width=\textwidth]{figures/fig1_architecture.png}
\caption{Architecture of the pharmacodynamically constrained conditional variational autoencoder (CVAE). Resting-state EEG features (2185-D) and ChemBERTa drug embeddings are encoded and fused with a disease label, then mapped by a CVAE whose training includes pharmacodynamic band and connectivity constraints. Training cohort: $N = 66$ (AD 37 / HC 29).}
\label{fig:architecture}
\end{figure*}


The EEG encoder maps structured tensors to a 128-D EEG latent. The drug encoder maps a frozen 384-D ChemBERTa embedding (\texttt{DeepChem/ChemBERTa-77M-MLM}) to a 64-D drug latent; ChemBERTa is not trained. Fusion concatenates EEG latent, drug latent, and a scalar disease label (AD = 1, HC = 0) into a 256-D representation. The CVAE encoder outputs a diagonal Gaussian in 128 dimensions with reparameterization
\begin{equation}
z = \mu + \exp(\tfrac{1}{2}\log\sigma^2)\odot\varepsilon, \qquad \varepsilon\sim\mathcal{N}(0,I).
\label{eq:reparam}
\end{equation}
The decoder reconstructs the 2185-D vector. At inference, baseline uses a zero drug vector; Donepezil/Memantine use their ChemBERTa embeddings. \(\mu_{\mathrm{base}}\) denotes the zero-drug latent mean.

Signed literature priors enter as penalty terms added to reconstruction and KL \citep{Babiloni2006DonepezilEEG,Jeong2004EEGreview,Winblad2007Memantine,Rive2013Memantine}:

\begin{itemize}
  \item Donepezil: \(\delta-0.15\), \(\theta-0.20\), \(\alpha+0.25\), \(\beta+0.10\), connectivity \(+0.20\)
  \item Memantine: \(\delta-0.10\), \(\theta-0.15\), \(\alpha+0.10\), \(\beta+0.05\), connectivity \(+0.15\)
\end{itemize}

With reconstruction \(L_{\mathrm{MSE}}=\mathrm{mean}((\hat{x}-x)^2)\), KL weight \(\beta=0.01\), band ReLU penalties \(L_{\mathrm{band}}\), connectivity non-decrease penalty \(L_{\mathrm{conn}}\) (prior connectivity magnitudes are stored but not multiplied into \(L_{\mathrm{conn}}\)), and schedule \(\alpha_c(t)\) rising from 0 after a 5-epoch warmup to \(\alpha_c^{\star}=0.1\), the total loss is
\begin{equation}
L = L_{\mathrm{MSE}} + \beta L_{\mathrm{KL}} + \alpha_c(t)\,L_{\mathrm{band}} + \alpha_c(t)\,L_{\mathrm{conn}}.
\label{eq:total}
\end{equation}
The KL term against \(\mathcal{N}(0,I)\) is
\begin{equation}
L_{\mathrm{KL}} = \frac{1}{B}\sum_{i=1}^{B}\Bigl(-\tfrac{1}{2}\sum_{k=1}^{128}\bigl(1+\log\sigma^2_{ik}-\mu_{ik}^2-\exp(\log\sigma^2_{ik})\bigr)\Bigr).
\label{eq:kl}
\end{equation}
For band \(b\in\{\delta,\theta,\alpha,\beta\}\) with reconstructed change \(\Delta_b\) versus baseline and prior \(\pi_b\),
\begin{equation}
L_{\mathrm{band}} = \sum_{b}|\pi_b|\,\mathrm{ReLU}\bigl(-\mathrm{sign}(\pi_b)\,\Delta_b\bigr),
\label{eq:band}
\end{equation}
and with connectivity-window change \(\Delta_c\),
\begin{equation}
L_{\mathrm{conn}} = \mathrm{ReLU}(-\Delta_c).
\label{eq:conn}
\end{equation}
The implemented schedule is
\begin{equation}
\alpha_c(t)=\begin{cases}
0 & t < T_w,\\
\alpha_c^{\star}\min\bigl(1,(t-T_w)/5\bigr) & t \ge T_w,
\end{cases}
\label{eq:alpha}
\end{equation}
with \(T_w=5\). Algorithm~\ref{alg:train} summarizes training.
\begin{algorithm}[H]
\caption{Training procedure for the pharmacodynamically constrained CVAE.}
\label{alg:train}
\small
\textbf{Input:} packed EEG features $x\in\mathbb{R}^{2185}$ (PSD, band powers, interleaved coherence/PLV); ChemBERTa embedding $d\in\mathbb{R}^{384}$; disease label $y\in\{0,1\}$; subject-level baseline features $x^{\mathrm{base}}$; priors $\pi$ from \texttt{get\_eeg\_effect\_prior}; $\beta=0.01$; $\alpha_c^{\star}=0.1$; $T_w=5$; Adam hyperparameters.\\
\textbf{Output:} checkpoint minimizing validation total loss.\\[0.4em]
\textbf{Training} (subject-level 5-fold; within-fold 10\% validation). For each epoch $t=0,\ldots,T-1$ and each minibatch:
\begin{enumerate}
  \setlength{\itemsep}{0.15em}
  \item Flatten EEG tensors to $x$ with the same packing as the decoder target.
  \item Encode EEG: $h_{\mathrm{eeg}}=\mathrm{EEGEncoder}(\mathrm{PSD},\mathrm{band},\mathrm{coh},\mathrm{PLV})$.
  \item Encode drug: $h_{\mathrm{drug}}=\mathrm{DrugEncoder}(d)$ (384-D ChemBERTa; enhanced 416-D prior concatenation is off for this variant).
  \item Fuse: $h=\mathrm{Fusion}(h_{\mathrm{eeg}},h_{\mathrm{drug}},y)$.
  \item CVAE encode: $(\mu,\log\sigma^2)=\mathrm{Enc}(h)$.
  \item Reparameterize: $z=\mu+\exp(\tfrac12\log\sigma^2)\odot\varepsilon$, $\varepsilon\sim\mathcal{N}(0,I)$.
  \item Decode: $\hat{x}=\mathrm{Dec}([z; h_{\mathrm{drug}}; y])$ (linear output, no output activation).
  \item $L_{\mathrm{MSE}}=\mathrm{mean}((\hat{x}-x)^2)$; $L_{\mathrm{KL}}$ as in Eq.~(\ref{eq:kl}).
  \item If the batch item is not baseline, compute $L_{\mathrm{band}}$ and $L_{\mathrm{conn}}$ vs $x^{\mathrm{base}}$ (Eqs.~(\ref{eq:band})--(\ref{eq:conn})); else both are 0.
  \item Set $\alpha_c(t)$ from Eq.~(\ref{eq:alpha}); form $L$ (Eq.~(\ref{eq:total})). Unique \texttt{drug\_id} subsets are scored separately and summed with batch-fraction weights.
  \item Backpropagate $L$ and take one Adam step.
\end{enumerate}
After each epoch, record train/val losses. Save the fold checkpoint with the lowest validation total loss.\\[0.3em]
\textbf{External evaluation} (not used for the Adam step): freeze the locked fold-0 checkpoint; simulate Donepezil/Memantine feature effects; score direction, magnitude, and encoding probes on held-out cohorts.
\end{algorithm}

Exact index slices and packing notes are in S1.

\subsection{Training cohort, freeze, and sensitivity protocol}

Training features come from resting-state EEG for 66 subjects (AD 37, HC 29; OpenNeuro ds004504). Preprocessing uses FIR 0.5-40 Hz, notch 50 Hz on the training path, FastICA with empty default exclusion, 4 s epochs at 90% overlap, and per-subject z-scoring before Welch/connectivity (S1). Features are 500 Hz spectral/connectivity packs. \(N=66\) is small relative to 2185-D inputs and multilayer capacity; primary claims are therefore external assessments, not training-set diagnosis.

Splits are subject-level stratified 5-fold (\texttt{random\_state=42}); within-fold validation is 10%. External cohorts are never in the training loop. Optimization uses Adam (\(10^{-4}\), weight decay \(10^{-5}\)), batch size 16, 100 epochs. Checkpoint selection minimizes internal validation total loss. The locked alias is fold-0 seed-42 epoch 99 (val total 0.0241). External \(r\), signs, and AUCs are post-freeze assessments, not selection criteria. The freeze is a documented project rule, not independently auditable preregistration. Training-seed sensitivity holds fold membership fixed (\texttt{split\_seed=42}) and varies \texttt{train\_seed} over \(\{42,7,21,123,2024\}\). Compute summary is in Table~\ref{tab:compute}.

\subsection{External cohorts}

Cohorts were introduced so each answered a doubt left by the previous setting (Table~\ref{tab:datasets}).
% Table 1 -- External assessment datasets (fit page width)
% Caption placement: ABOVE the tabular (Elsevier convention).
% Requires: booktabs, adjustbox
\begin{table*}[!t]
\centering
\footnotesize
\setlength{\tabcolsep}{3pt}
\caption{External EEG datasets used for escalating external assessment of the constrained twin.
Sampling rates list native acquisition then the common 500\,Hz analysis rate.
Roles follow the Results narrative sequence.}
\label{tab:datasets}
\begin{adjustbox}{max width=\textwidth}
\begin{tabular}{@{}llp{2.2cm}p{2.4cm}p{2.5cm}p{1.9cm}p{2.4cm}@{}}
\toprule
Dataset & N & Groups & Labels & Sampling rate & Montage & Role \\
\midrule
TUH & 200 & Abn 131 / Nor 69 & Abnormal vs normal & Mixed EDF $\rightarrow$ 500\,Hz & 19ch 10--20 & First external; proxy labels \\
OSF & 92 & AD 80 / HC 12 & Folder AD vs HC & 256 $\rightarrow$ 500\,Hz & 19ch 10--20 & AD/HC; small HC \\
P-ADIC & 145 & AD 49 / HC 96 & NIA-AA AD vs HC & 200/500 $\rightarrow$ 500\,Hz & 19ch 10--20 & Larger AD/HC; null check \\
CAUEEG & 1122 & Nor 436 / MCI 395 / Dem 291 & Normal / MCI / Dementia & 200 $\rightarrow$ 500\,Hz & 19ch AVG 10--20 & Large-N; encoding \\
\bottomrule
\end{tabular}
\end{adjustbox}
\end{table*}


\begin{itemize}
  \item \textbf{TUH} (\(n=200\); abnormal 131 / normal 69): locked features use \texttt{crop\_seconds=60.0} and \texttt{apply\_ica=false} in \texttt{processing\_report.json}; current script default \texttt{CROP\_SECONDS=320.0} was not used to regenerate locked results.
  \item \textbf{OSF} AD/HC (\(n=92\); HC \(n=12\)).
  \item \textbf{P-ADIC} AD/HC (\(n=145\)) \citep{Shor2021PADIC}.
  \item \textbf{CAUEEG} Normal/MCI/Dementia (1122 features OK) for magnitude and encoding \citep{Kim2023CAUEEG}.
\end{itemize}
Direction tests use TUH, OSF, and P-ADIC. CAUEEG was not required for the locked 10/10 direction claim.

\subsection{Evaluation endpoints and statistics}

No paired pre/post-dose EEG is available. Direction measures concordance with the \textbf{training constrained signature}, a prior-shaped model-derived reference. That circularity is intentional and bounded by controls: prior-only (constant literature-signed effect; mean \(r=0.636\), 10/10), matched unconstrained CVAE (\(\alpha_c=0\)), and a ridge imitation control trained on locked CVAE-simulated targets (not observed post-dose EEG).

\textbf{Endpoint families (SEME).} (1) \textbf{S / E:} Effect-vector Pearson \(r\) (primary continuous score) and block-level signed agreement (10/10 secondary) on cohort-mean drug-minus-baseline vectors. (2) \textbf{M:} Magnitude \(\bar{s}=\tfrac12(\|\mu_{\mathrm{don}}-\mu_{\mathrm{base}}\|_2+\|\mu_{\mathrm{mem}}-\mu_{\mathrm{base}}\|_2)\). (3) \textbf{E (encoding):} CAUEEG probes on \(\mu_{\mathrm{base}}\) (disease bit fixed to 0) versus theta/alpha and packed 2185-D logistic references; nested heads use outer 5-fold / inner 3-fold search. (4) Exploratory \(\alpha_c\) sweep on fold-0 twins. (5) Paired DeLong tests on stored subject-level OOF probe scores (\(n=727\)) \citep{DeLong1988}. (6) Matched fold-0 ChemBERTa-versus-padded-one-hot ablation (same split, seed 42, \(\alpha_c=0.1\), DrugEncoder input dim 384; non-locked checkpoints only). Uncertainty for \(r\) uses subject-bootstrap percentile intervals (B = 2000). Magnitude uses Mann-Whitney and Cohen's d. Probe permutation tests use \(N_{\mathrm{perm}}=200\). No multiplicity correction was applied across the escalation sequence or DeLong pairs. Full endpoint definitions are in S1.
